\section{Point de d\'epart et cible}\label{sec:depart}
% Chapitre 2 (tache 10, passage chapitres 02-03). Figures: fig_ziel_nominal,
% fig_kaskade (voir docs/figures.md). Aucun chiffre tape: macros de
% data/kennzahlen.tex, generees par scripts/rechnung_retraite.py.
% Titres d'axes et etiquettes de barres passes par l'argument optionnel de
% \linienfigur/\balkenfigur (correctif infrastructure, tache 10, second
% passage): remplace le contournement "every axis post" du premier passage.

Ce chapitre r\'epond \`a deux questions pr\'ealables \`a tout calcul~: d'o\`u pars-tu, et combien de
dividendes bruts faut-il encaisser pour qu'il te reste \ZielMonat~€ par mois \`a d\'epenser~? La
section~\ref{sec:depart-profil} fixe ton point de d\'epart, la section~\ref{sec:depart-cible} exprime
la cible en euros constants puis en euros nominaux, et la section~\ref{sec:depart-cascade} descend du
dividende brut au montant disponible.

\subsection{Ton point de d\'epart}\label{sec:depart-profil}

Tu es n\'e en 2001, tu termines ton master en ao\^ut 2027 apr\`es ton stage chez Jungheinrich, et tu
vis en Allemagne, o\`u tu resteras r\'esident fiscal pendant la vie active comme \`a la retraite. Tu es
c\'elibataire, sans enfant et sans bien immobilier, et tu acceptes un risque \'elev\'e sur tes
placements. Ton \'epargne commence en septembre 2027, avec au moins \EpargneMinimum~€ par mois quoi
qu'il arrive, puis un taux d'\'epargne de \SparquoteZehn{}, \SparquoteZwanzig{}, \SparquoteDreissig{} ou
\SparquoteFuenfzig~\% de ton salaire net~; le chapitre~\ref{sec:accumulation} construit ce salaire \`a partir d'une
r\'ef\'erence publi\'ee et calcule l'\^age de d\'epart qui en r\'esulte.

Ton capital de d\'epart exact n'est pas connu \`a ce jour. Le plan travaille donc sur une fourchette
de \KapitalDepartNiedrig~€ \`a \KapitalDepartHoch~€ et retient son milieu, \KapitalDepartMitte~€, pour
le sc\'enario de base. Cette incertitude p\`ese peu face au capital \`a r\'eunir~: le chapitre~\ref{sec:capital} chiffre
\`a \KapitalEtfReferenz~€ celui de la seule strat\'egie de r\'ef\'erence en ETF, et le chapitre~\ref{sec:accumulation}
montre que les deux bornes de la fourchette donnent le m\^eme \^age de d\'epart. Les
startups, le trading intraday et les paris sur des titres isol\'es restent hors du calcul, comme le
rappelle l'annexe (section~\ref{sec:annexe-perimetre}).

Le sc\'enario de base, auquel renvoient tous les encadr\'es \enquote{Yann}, combine ce capital de
d\'epart, la r\'epartition \ReferenzStrategie{} entre ETF maison et ETF monde, et une \'epargne de
\ReferenzSparquote~\% du salaire net. Ce taux n'a pas \'et\'e choisi par confort~: \`a
\SparquoteZwanzig~\%, la m\^eme strat\'egie de dividendes n'atteint jamais l'objectif dans l'horizon de
calcul, alors qu'un ETF monde avec un retrait de \TauxRetrait~\% par an y parvient \`a
\AlterEtfReferenzZwanzig~ans. Le chapitre~\ref{sec:accumulation} d\'etaille ce r\'esultat, le plus d\'efavorable du
document pour la th\`ese des dividendes.

\subsection{La cible, en euros constants puis nominaux}\label{sec:depart-cible}

L'objectif est de disposer de \ZielMonat~€ par mois, soit \ZielJahr~€ par an, en pouvoir d'achat de
2026, apr\`es imp\^ot et apr\`es cotisations d'assurance maladie et d\'ependance. Tout le document
raisonne en euros de 2026, que l'on appelle euros constants~: un montant de \ZielMonat~€ en 2050
d\'esigne ce que \ZielMonat~€ ach\`etent aujourd'hui, et non la somme inscrite ce jour-l\`a sur ton
relev\'e. Ce choix rend comparables des montants s\'epar\'es par plusieurs d\'ecennies, mais il exige
deux hypoth\`eses explicites, l'une sur l'inflation et l'autre sur le rendement r\'eel du march\'e.

\annahme{Inflation de \InflationZiel~\% par an sur toute la dur\'ee du plan, soit l'objectif de
stabilit\'e des prix de la Banque centrale europ\'eenne\QH{inflation_ziel}. Les variantes \`a
\InflationBas~\% et \InflationHaut~\% de la figure~\ref{fig:cible-euros-nominaux} servent seulement
\`a mesurer la sensibilit\'e de la conversion en euros nominaux.}

\annahme{Rendement r\'eel du march\'e d'actions de \RenditeReal~\% par an, calcul\'e sur les
donn\'ees de Robert Shiller pour l'indice \IndiceActionsUs\QH{shiller_url}~: \RenditeReelleConvention. Le plan
initial pr\'evoyait la m\'ediane des rendements r\'eels annuels, qui vaut \RenditeReelleMediane~\%
sur la m\^eme p\'eriode. Cet \'ecart au plan est assum\'e~: une m\'ediane de rendements annuels
n'est pas un taux de capitalisation, car elle ignore l'effet multiplicatif des ann\'ees de krach, et
l'appliquer sur plusieurs d\'ecennies surestimerait fortement le capital accumul\'e.}

Pour lire un article de presse ou un relev\'e de courtier dans vingt ans, il faudra faire le chemin
inverse et convertir les euros nominaux de l'\'epoque en euros de 2026. La
figure~\ref{fig:cible-euros-nominaux} montre ce que devient la cible mensuelle sous trois hypoth\`eses
d'inflation.

\linienfigur{fig_ziel_nominal}[xlabel={Ann\'ee}, ylabel={Cible mensuelle (euros nominaux)},
  xticklabel style={/pgf/number format/set thousands separator={}}, scaled y ticks=false,
  legend pos=north west]%
  {jahr}{inflation_1,inflation_2,inflation_3}%
  {{Inflation \InflationBas~\%},{Inflation \InflationZiel~\%},{Inflation \InflationHaut~\%}}%
  {Cible de \ZielMonat~€ par mois en euros de 2026, convertie en euros nominaux de chaque ann\'ee
  sous trois hypoth\`eses d'inflation.\label{fig:cible-euros-nominaux}}

Ce graphique montre qu'\`a l'inflation cible de \InflationZiel~\%, les \ZielMonat~€ d'aujourd'hui
correspondent \`a \ZielNominalXLIV~€ par mois en 2044 et \`a \ZielNominalLX~€ en 2060. Si l'inflation
tenait \InflationHaut~\% par an, il faudrait \ZielNominalLXHaut~€ nominaux en 2060, et
\ZielNominalLXBas~€ seulement \`a \InflationBas~\%. L'\'ecart entre les deux courbes extr\^emes montre
qu'un objectif fix\'e en euros nominaux, sans r\'ef\'erence \`a une ann\'ee de prix, ne veut rien
dire \`a cet horizon.

\subsection{Du dividende brut au montant disponible}\label{sec:depart-cascade}

Un dividende brut n'arrive jamais entier sur ton compte. L'\'Etat de la soci\'et\'e peut en retenir
une partie \`a la source, le fisc allemand pr\'el\`eve ensuite son imp\^ot, et la caisse d'assurance
maladie calcule ses cotisations sur le dividende brut. Le chapitre~\ref{sec:fiscalite} d\'etaille
chacune de ces \'etapes~; il suffit ici d'en voir l'ordre de grandeur.

\annahme{Cascade calcul\'ee avec les r\`egles fiscales et sociales de 2026, en euros de 2026, pour un
revenu venant uniquement d'actions d\'etenues en direct, avec la r\'epartition par pays du
portefeuille-exemple du chapitre~\ref{sec:portefeuille}~: \PartUs~\% de dividendes am\'ericains, \PartDe~\% allemands et
\PartFr~\% fran\c{c}ais. \HypotheseUsFormulaire. Imp\^ot au taux forfaitaire, sans la
G\"unstigerpr\"ufung discut\'ee au chapitre~\ref{sec:fiscalite}.}

Sous ces hypoth\`eses, il faut \BruttoNoetig~€ de dividendes bruts par an, soit \BruttoNoetigMois~€
par mois, pour qu'il reste \ZielJahr~€ disponibles. La figure~\ref{fig:cascade-brut-disponible}
d\'ecompose le passage de l'un \`a l'autre.

% Ordre des etiquettes = ordre des lignes de data/fig_kaskade.csv (brut_necessaire,
% retenues_etrangeres, impot_allemand, kv_pv, disponible): la colonne "posten" porte
% des "_" (compile desormais sans risque via "posten_disp", voir preamble.tex), mais
% des libelles francais restent plus lisibles, d'ou la liste explicite ci-dessous.
\balkenfigur{fig_kaskade}[ylabel={Montant annuel (euros de 2026)},
  scaled y ticks=false, bar width=22pt, enlarge y limits=0.25,
  xticklabels={{Dividende brut},{Retenues \'etrang\`eres},{Imp\^ot allemand},{Maladie et d\'ependance},{Disponible}},
  nodes near coords, every node near coord/.append style={font=\scriptsize,
    /pgf/number format/fixed, /pgf/number format/precision=0}]%
  {posten}{betrag}%
  {Du dividende brut n\'ecessaire au montant disponible, sur une ann\'ee de retraite (r\`egles de
  2026, euros de 2026, actions en direct).\label{fig:cascade-brut-disponible}}

Ce graphique montre que sur \BruttoNoetig~€ de dividendes bruts, il ne reste que \ZielJahr~€, soit
\PartDisponible~\% du brut. Les retenues \'etrang\`eres prennent \RetenuesEtrangeres~€, l'imp\^ot
allemand \ImpotAllemand~€, et l'assurance maladie et d\'ependance \KvPvAnnuel~€, soit
\KvPvAnteilBrut~\% du brut~: ce dernier poste, rarement cit\'e dans les articles sur la rente par
dividendes, p\`ese plus lourd que l'imp\^ot allemand lui-m\^eme. Il serait encore plus \'elev\'e sans
le plafond de cotisation, que ton revenu brut mensuel d\'epasse d\'ej\`a.

\bk{Retiens de ce chapitre qu'il faut raisonner en brut pour dimensionner le capital, et en euros de
2026 pour comparer des ann\'ees entre elles~; un calcul fait sur le net ou en euros nominaux sous-estime
le capital n\'ecessaire.} Le chapitre~\ref{sec:fiscalite} explique d'o\`u vient chaque poste de cette
cascade, et ce qui change quand le dividende vient d'un autre pays ou d'un ETF.

\Yannbox{\AnneeYannBasis}{Dans le sc\'enario de base (\KapitalDepartMitte~€ de d\'epart, r\'epartition
\ReferenzStrategie{} entre ETF maison et ETF monde, \ReferenzSparquote~\% d'\'epargne), Yann atteint son
objectif \`a \AlterYannBasis~ans, l'\^age l\'egal de la retraite~: il ne part pas plus t\^ot que la
loi ne le permet d\'ej\`a. Ses \ZielMonat~€ de 2026 valent alors \ZielNominalBasis~€ nominaux par
mois. Son capital de \KapitalYannBasis~€, exprim\'e en euros de 2026, lui verse \RevenuNetYannBasis~€
nets par mois, toujours en euros de 2026.}
